\begin{abstract}
Speech and language models need a written form for pronunciation that a tokenizer can handle well. The International Phonetic Alphabet is accurate and well suited to its own purpose, but less than ideal in this setting. A single sound is often a base letter followed by several combining marks, some symbols sit in the supplementary planes of Unicode, and a subword or byte tokenizer splits one pronunciation across several tokens with no fixed relation to the sounds. Talk is a phonetic encoding built for this setting. It writes every sound as a base consonant or vowel plus a small set of modifiers, in plain ASCII that a normal keyboard can type. It also maps each canonical sound to a single code point, so a whole word becomes one token per sound, using the Hangul syllable block for individual sounds and Chinese characters for consonant and vowel clusters. The map from sound to code point is fixed and append-only, so a model built on it keeps working as the inventory grows. Talk converts to and from IPA, without loss in the talk direction, and it splits words into syllables. It has no runtime dependencies and ships for Python and TypeScript. This paper describes the encoding, the token form, and how to use it. The system is implemented and in use, and coverage of the remaining IPA symbols is being completed.
\end{abstract}
